\documentclass[11pt]{article}

\usepackage[backend=bibtex,natbib,style=numeric-comp,sorting=none,doi=false,isbn=false,url=false]{biblatex}
\addbibresource{refs}

\usepackage[T1]{fontenc}
\usepackage{lmodern}
\usepackage{amsmath,amssymb}
\usepackage{hyperref}
\usepackage[margin=1in]{geometry}

\newcommand{\Aeff}{\alpha'_{\mathrm{eff}}}
\newcommand{\Sym}{\mathrm{Sym}}
\newcommand{\mc}[1]{\mathcal{#1}}
\newcommand{\bC}{\mathbb{C}}
\newcommand{\bR}{\mathbb{R}}
\newcommand{\ie}{\begin{equation}\begin{aligned}}
\newcommand{\fe}{\end{aligned}\end{equation}}

\title{A Curved First-Order String for \texorpdfstring{$\Sym^N(\bR^8)$}{SymN(R8)} on a Two-Dimensional Spacetime}
\author{}
\date{}

\begin{document}
\maketitle

\section{Aim}
This note records a concrete conjectural starting point for a string dual of the symmetric-orbifold SCFT
\begin{equation}\label{eq:proposed-dual-cft}
\Sym^N(\bR^8)
\end{equation}
living on a two-dimensional spacetime $C$. More explicitly, this means the $S_N$ orbifold of $N$ copies of the free $(8,8)$ SCFT of eight noncompact bosons and fermions. Here $C$ is itself the target of the longitudinal first-order string variables. The intended picture is close in spirit to matrix string theory
\cite{dijkgraaf1997matrixstring,arutyunov1998orbifoldscattering,dijkgraaf2003matrixstringcontactterms}. It is also close to the recent first-order-string approach to two-dimensional Yang-Mills \cite{komatsu2025stringduals2dym}, but with the two-dimensional maximally supersymmetric gauge theory as the field-theory side.

The status is mixed. The local worldsheet system is explicit and is just the supersymmetric first-order string already present in the Gomis--Ooguri description of the NCOS/NRCS limit \cite{gomis2000nonrelativisticclosedstringtheory}. The extension from flat target $C=\bC$ to a general Riemann surface $C$ is also explicit at the level of local patching data. What is \emph{not} yet proved is the full global BRST construction on arbitrary $C$, the precise joining/splitting interaction, and the quantum statement that physical observables depend only on the conformal structure of $C$ rather than on a chosen metric representative. Those points are part of the program, not finished results.

Throughout, $\Sigma$ denotes the worldsheet Riemann surface, $C$ denotes the two-dimensional target/spacetime, and $\Aeff$ denotes the effective string scale inherited from the non-relativistic string limit. Because the transverse theory is noncompact, any partition function involving the $\bR^8$ seed theory is understood either per unit transverse volume or, equivalently for the free formulas used below, after factoring out the center-of-mass zero modes. Whenever a formula is imported from \cite{gomis2000nonrelativisticclosedstringtheory} or from the curved-$\beta\gamma$ literature \cite{nekrasov2005curvedbetagammasystems}, I will say so explicitly.

\section{Flat-Space Starting Point}
The local matter system is taken as input. Concretely, I start from the critical type-II first-order string built from a longitudinal super-$\beta\gamma$ multiplet together with eight transverse free superfields, namely the local Gomis--Ooguri matter system written in $(1,1)$ superspace. On a Euclidean worldsheet with coordinates $(z,\bar z,\vartheta,\bar\vartheta)$, define
\begin{equation}\label{eq:superderivatives}
D=\partial_\vartheta+\vartheta\,\partial,
\qquad
\bar D=\partial_{\bar\vartheta}+\bar\vartheta\,\bar\partial,
\qquad
D^2=\partial,
\qquad
\bar D^2=\bar\partial.
\end{equation}
The flat-space action is
\begin{equation}\label{eq:flat-action}
\begin{aligned}
S_{\mathrm{flat}}
&=
\frac{1}{2\pi}\int d^2z\,d\vartheta\,d\bar\vartheta\,
\left[
\mathbf B\,\bar D\mathbf X^+
+\bar{\mathbf B}\,D\mathbf X^-
+\frac{1}{4\Aeff}\,D\mathbf X^+\,\bar D\mathbf X^-
\right]
+S_\perp,
\\
S_\perp
&=
\frac{1}{2\pi \Aeff}\int d^2z\,d\vartheta\,d\bar\vartheta\,
\delta_{ij}\,D\mathbf Y^i\,\bar D\mathbf Y^j,
\qquad i=1,\dots,8.
\end{aligned}
\end{equation}
Here $\mathbf X^\pm$ are bosonic longitudinal superfields, $\mathbf B$ and $\bar{\mathbf B}$ are fermionic first-order superfields, and $\mathbf Y^i$ with $i=1,\dots,8$ are transverse free superfields.
In Lorentzian signature one may think of $X^\pm=X^0\pm X^1$. For the curved-target discussion below it is more convenient to view $X^+$ as a local holomorphic coordinate on the target Riemann surface $C$ and $X^-$ as the corresponding antiholomorphic coordinate.

In Euclidean signature the physical contour should impose the reality condition
\begin{equation}\label{eq:euclidean-reality}
X^-=\overline{X^+},
\end{equation}
even though the first-order path integral is most conveniently written in terms of independent fields $\mathbf X^\pm$ before one chooses that contour.

Varying $\mathbf B$ and $\bar{\mathbf B}$ in \eqref{eq:flat-action} gives the first-order bulk constraints
\begin{equation}\label{eq:bulk-chiral-constraints}
\bar D\mathbf X^+=0,
\qquad
D\mathbf X^-=0,
\qquad
\bar D\mathbf B=0,
\qquad
D\bar{\mathbf B}=0.
\end{equation}
So on shell the longitudinal superfields are chiral,
\begin{equation}\label{eq:onshell-plus}
\mathbf X^+(z,\vartheta)=\gamma(z)+\vartheta\,c(z),
\qquad
\mathbf B(z,\vartheta)=b(z)+\vartheta\,\beta(z),
\end{equation}
\begin{equation}\label{eq:onshell-minus}
\mathbf X^-(\bar z,\bar\vartheta)=\bar\gamma(\bar z)+\bar\vartheta\,\bar c(\bar z),
\qquad
\bar{\mathbf B}(\bar z,\bar\vartheta)=\bar b(\bar z)+\bar\vartheta\,\bar\beta(\bar z),
\end{equation}
with operator products
\begin{equation}\label{eq:ope-holomorphic}
\beta(z)\gamma(w)\sim \frac{1}{z-w},
\qquad
b(z)c(w)\sim \frac{1}{z-w},
\end{equation}
\begin{equation}\label{eq:ope-antiholomorphic}
\bar\beta(\bar z)\bar\gamma(\bar w)\sim \frac{1}{\bar z-\bar w},
\qquad
\bar b(\bar z)\bar c(\bar w)\sim \frac{1}{\bar z-\bar w}.
\end{equation}
The component conformal weights are
\begin{equation}\label{eq:component-weights}
h(\gamma)=0,
\qquad
h(\beta)=1,
\qquad
h(c)=h(b)=\frac{1}{2},
\end{equation}
and similarly in the antiholomorphic sector. Hence the bosonic $(\beta,\gamma)$ system contributes $c=2$, while the fermionic $(b,c)$ system contributes $c=1$, so
\begin{equation}\label{eq:longitudinal-central-charge}
c_{+-}=2+1=3.
\end{equation}
while the eight transverse bosons and fermions contribute $12$, so
\begin{equation}\label{eq:total-central-charge}
c_{\mathrm{matter}}=15.
\end{equation}
Thus the matter system is exactly that of a critical type-II RNS string, written in first-order variables. The first two terms in \eqref{eq:flat-action} are the Poisson/BF core, while the $1/\Aeff$ term keeps a finite string scale rather than collapsing the theory to a purely topological model.

It is useful to say explicitly what the finite $\Aeff$ deformation does and does not do. Because the equations of motion \eqref{eq:bulk-chiral-constraints} come from varying $\mathbf B$ and $\bar{\mathbf B}$, the last term in \eqref{eq:flat-action} does not change the classical chirality constraints. Nevertheless it is not physically trivial: it fixes the normalization of global energies and chemical potentials, and the same $\Aeff$ reappears later in the genus-one energy formula \eqref{eq:go-energy}. So the deformation is on-shell trivial only at the level of the local bulk equations, not at the level of observables. Its coefficient is also not a free choice: the relative normalization of the BF term and the $1/\Aeff$ deformation is fixed by the original Gomis--Ooguri scaling limit.

\section{Patchwise Definition on a Riemann Surface}
Let $C$ be a target Riemann surface. Choose holomorphic charts $U_a\subset C$ with local coordinate $u_a$ and local metric representative
\begin{equation}\label{eq:local-target-metric}
ds_C^2=e^{2\omega_a(u_a,\bar u_a)}\,du_a\,d\bar u_a.
\end{equation}
On the worldsheet we introduce, in each chart, superfields
\begin{equation}\label{eq:local-target-fields}
\mathbf X^+_a,\quad \mathbf X^-_a,\quad \mathbf B_a,\quad \bar{\mathbf B}_a
\end{equation}
and define the local longitudinal action
\begin{equation}\label{eq:curved-local-action}
\begin{aligned}
S^{(a)}_{\mathrm{long}}
=
\frac{1}{2\pi}\int d^2z\,d\vartheta\,d\bar\vartheta\,
\left[
\mathbf B_a\,\bar D\mathbf X^+_a
+\bar{\mathbf B}_a\,D\mathbf X^-_a
+\frac{e^{2\omega_a(\mathbf X_a^+,\mathbf X_a^-)}}{4\Aeff}\,
D\mathbf X^+_a\,\bar D\mathbf X^-_a
\right].
\end{aligned}
\end{equation}
The transverse action $S_\perp$ is unchanged, since the transverse target remains $\bR^8$.

This is precisely the place where the curved $\beta\gamma$ machinery reviewed by Nekrasov \cite{nekrasov2005curvedbetagammasystems} is essential. The first-order fields do not patch as ordinary scalars; they patch as a coordinate and its conjugate one-form under target holomorphic coordinate changes. Without this machinery one only has a local action in each chart, not a globally defined worldsheet theory on $C$.

On an overlap $U_a\cap U_b$, let
\begin{equation}\label{eq:target-transition-functions}
u_b=f_{ba}(u_a),
\qquad
\bar u_b=\bar f_{ba}(\bar u_a),
\end{equation}
with $f_{ba}$ holomorphic. The superfields are patched by
\begin{equation}\label{eq:superfield-patching}
\mathbf X^+_b=f_{ba}(\mathbf X^+_a),
\qquad
\mathbf X^-_b=\bar f_{ba}(\mathbf X^-_a),
\end{equation}
\begin{equation}\label{eq:beta-patching}
\mathbf B_b=\big(f'_{ba}(\mathbf X^+_a)\big)^{-1}\mathbf B_a,
\qquad
\bar{\mathbf B}_b=\big(\bar f'_{ba}(\mathbf X^-_a)\big)^{-1}\bar{\mathbf B}_a.
\end{equation}
This is the natural superfield version of a curved $\beta\gamma$ system. Indeed,
\begin{equation}\label{eq:bf-patching-check}
\mathbf B_b\,\bar D\mathbf X^+_b
=
\big(f'_{ba}(\mathbf X^+_a)\big)^{-1}\mathbf B_a\,
f'_{ba}(\mathbf X^+_a)\bar D\mathbf X^+_a
=
\mathbf B_a\,\bar D\mathbf X^+_a,
\end{equation}
and similarly for the antiholomorphic term. The metric factor transforms as
\begin{equation}\label{eq:metric-transformation}
e^{2\omega_b(u_b,\bar u_b)}
=
e^{2\omega_a(u_a,\bar u_a)}
\left|\frac{du_a}{du_b}\right|^2,
\end{equation}
so the last term in $S^{(a)}_{\mathrm{long}}$ also patches correctly:
\begin{equation}\label{eq:kinetic-patching-check}
e^{2\omega_b}D\mathbf X^+_b\,\bar D\mathbf X^-_b
=
e^{2\omega_a}D\mathbf X^+_a\,\bar D\mathbf X^-_a.
\end{equation}
The chirality constraints are also coordinate-independent:
\begin{equation}\label{eq:chirality-patching}
\bar D\mathbf X^+_b=f'_{ba}(\mathbf X^+_a)\bar D\mathbf X^+_a,
\qquad
D\mathbf X^-_b=\bar f'_{ba}(\mathbf X^-_a)D\mathbf X^-_a.
\end{equation}
So the statement that the bosonic map $\gamma:\Sigma\to C$ is holomorphic is globally meaningful and does not depend on the choice of target chart.
The cocycle condition on triple overlaps follows automatically from
\begin{equation}\label{eq:cocycle}
f_{ca}=f_{cb}\circ f_{ba}
\end{equation}
together with the chain rule. Moreover, the Euclidean reality condition \eqref{eq:euclidean-reality} is compatible with \eqref{eq:superfield-patching} because for a genuine real two-manifold $C$ the antiholomorphic transition function is the complex conjugate $\bar f_{ba}$ of $f_{ba}$. If $\gamma:\Sigma\to C$ denotes the bosonic map underlying the superfield $\mathbf X^+$, then the component fields $c$ and $b$ in \eqref{eq:onshell-plus} are naturally sections of $\gamma^*(TC)$ and $\gamma^*(T^*C)$.

It is also useful to be explicit about the fermions. The fermionic components $c$ and $\bar c$ transform as sections of $\gamma^*(TC)$ and $\gamma^*(\overline{TC})$, while $b$ and $\bar b$ transform as sections of the dual bundles. They are not target-space spinors. So the patching data above do not require an independent spin structure on $C$; the genuine spin-structure subtlety appears later when $C$ itself is a spacetime torus and one chooses periodic or antiperiodic boundary conditions for spacetime fermions around its cycles.

One can also check that the curved deformation does not spoil the on-shell chirality of $\mathbf B$. Varying $\mathbf X^+_a$ gives
\begin{equation}\label{eq:curved-b-eom}
\bar D\mathbf B_a
=
\frac{1}{4\Aeff}
\left[
\partial_{+}\!\big(e^{2\omega_a}\big)\,D\mathbf X^+_a\,\bar D\mathbf X^-_a
-D\!\left(e^{2\omega_a}\bar D\mathbf X^-_a\right)
\right],
\end{equation}
where $\partial_+$ denotes differentiation with respect to the holomorphic target coordinate $u_a$. Using the first-order constraints $\bar D\mathbf X^+_a=0$ and $D\mathbf X^-_a=0$, together with $D\bar D\mathbf X^-_a=-\bar DD\mathbf X^-_a=0$, the right-hand side of \eqref{eq:curved-b-eom} vanishes on shell. Hence $\bar D\mathbf B_a=0$, and similarly $D\bar{\mathbf B}_a=0$, also hold in the curved theory.

Therefore the local first-order theory defines a globally consistent \emph{classical} worldsheet action on any target Riemann surface $C$. This classical patching statement is derived in the note; it is not merely conjectural. Its classical solution locus is the moduli space of holomorphic maps $\gamma:\Sigma\to C$. For $C=T^2$ this locus reduces to the familiar discrete sum over winding data, whereas for general $C$ it is a genuine moduli space of holomorphic maps. Any extension of the torus free-energy check away from $C=T^2$ would therefore be expected to replace the discrete $(w,n)$ sum by an integral, or at least a sum over connected components, of this holomorphic-map moduli space together with the corresponding fluctuation determinants.

\section{Quantum Status and the Target-Space Weyl Question}
The first-order Poisson/BF core,
\begin{equation}\label{eq:bf-core}
\frac{1}{2\pi}\int d^2z\,d\vartheta\,d\bar\vartheta\,
\left(
\mathbf B\,\bar D\mathbf X^+
+\bar{\mathbf B}\,D\mathbf X^-
\right),
\end{equation}
does not involve the target metric at all. The only place where the chosen local representative
\begin{equation}\label{eq:target-weyl-representative}
ds_C^2=e^{2\omega}du\,d\bar u
\end{equation}
enters is the finite string-tension deformation
\begin{equation}\label{eq:metric-deformation}
\frac{1}{8\pi \Aeff}\int d^2z\,d\vartheta\,d\bar\vartheta\,
e^{2\omega(\mathbf X^+,\mathbf X^-)}D\mathbf X^+\,\bar D\mathbf X^-.
\end{equation}
So classically the dependence on a Weyl representative is controlled: it affects only the term \eqref{eq:metric-deformation} and does not alter the holomorphic patching of the first-order algebra. Quantum mechanically, however, this is not enough to conclude target-space Weyl invariance.

There are two distinct issues. First, the composite operator $e^{2\omega(\mathbf X^+,\mathbf X^-)}$ in \eqref{eq:metric-deformation} requires a regularization and normal-ordering prescription. Because the bulk $X^+X^-$ OPE is regular in the first-order variables, the short-distance singularities are milder than in an ordinary second-order sigma model, but finite local counterterms can still renormalize the curved deformation.

The mechanism behind this improved UV behavior is simple. In the strict BF limit obtained by dropping the $1/\Aeff$ term, the only singular OPEs are $\beta(z)\gamma(w)\sim 1/(z-w)$, $b(z)c(w)\sim 1/(z-w)$, and their antiholomorphic counterparts; the mixed OPE $X^+(z,\bar z)X^-(w,\bar w)$ is regular. So the composite operator $e^{2\omega(X^+,X^-)}$ does not already carry the standard second-order normal-ordering singularity. Such singularities arise only after the finite-$\Aeff$ deformation is turned on, which is why the quantum problem is milder than in the usual sigma-model presentation, even though it is not absent.

Second, one expects, at minimum, the usual curved-$\beta\gamma$ measure anomaly of the type emphasized in \cite{nekrasov2005curvedbetagammasystems}. For a target Riemann surface $C$, its topological part is controlled by the first Chern class of $TC$, hence by
\begin{equation}\label{eq:target-euler}
\int_C c_1(TC)=\chi(C)=2-2g_C.
\end{equation}
Here $g_C$ is the genus of the target Riemann surface $C$. So for $C=S^2$ the topological obstruction is certainly nontrivial, for $C=T^2$ it vanishes, and for $g_C>1$ it is again nonzero. This does not by itself determine the full local counterterm problem, but it does show that target-space Weyl independence is not automatic. In a full string construction this may require a local counterterm or dilaton-like worldsheet curvature coupling of the form
\begin{equation}\label{eq:dilaton-like-coupling}
S_\Phi=\frac{1}{4\pi}\int_\Sigma \sqrt{g_{\rm ws}}\,R^{(2)}_{\rm ws}\,\Phi(X^+,X^-).
\end{equation}
Here $g_{\rm ws}$ is a chosen metric on the worldsheet $\Sigma$, $R^{(2)}_{\rm ws}$ is its scalar curvature, and $\Phi$ is a possible target-space scalar counterterm. The present note does not determine that term.

A full quantum construction would also require a renormalized worldsheet stress tensor and supercurrent on curved $C$. In flat space they are those of the standard first-order RNS system and give the critical matter central charge \eqref{eq:total-central-charge}. On curved $C$, however, the deformation \eqref{eq:metric-deformation} generates $\omega$-dependent composite corrections to both operators. Showing that these corrected operators still obey the appropriate OPEs, preserve $c_{\rm matter}=15$, and define the BRST operator globally is part of the open quantum completion problem.

This is the right place to distinguish two different notions:
\begin{itemize}
\item Worldsheet conformal invariance is the statement that the local matter system has the critical central charge \eqref{eq:total-central-charge}.
\item Target-space Weyl invariance is the much stronger statement that changing
\begin{equation}\label{eq:target-weyl-rescaling}
\omega\longmapsto \omega+\sigma
\end{equation}
does not change the physical BRST data after all required counterterms are included.
\end{itemize}
The current note only establishes the first statement. The second is a conjectural quantum completion problem. The honest claim is therefore narrower than in the earlier draft: curved $\beta\gamma$ patching under holomorphic coordinate changes is explicit and essential, whereas target-space Weyl independence remains to be proved after the anomaly problem and the normal ordering of \eqref{eq:metric-deformation} are understood.

\section[Conjectural Duality to SymN(R8)]{Conjectural Duality to $\Sym^N(\bR^8)$}
The proposed dual field theory is the symmetric-orbifold SCFT
\begin{equation}\label{eq:dual-orbifold}
\Sym^N(\bR^8)
\end{equation}
placed on the same two-dimensional spacetime $C$. Concretely, this is the $S_N$ orbifold of $N$ copies of the free $(8,8)$ $\bR^8$ sigma model. The duality claim itself is conjectural. Equivalently, one is proposing that the same theory is the infrared limit of two-dimensional $(8,8)$ $U(N)$ SYM in the standard matrix-string regime
\cite{dijkgraaf1997matrixstring,cho2026fluxsectorsmatrixstring}.
For $C=S^1\times \bR$ and for $C=T^2$, this statement is directly motivated by the usual matrix-string logic. For a general Riemann surface $C$, it should presently be regarded as an extrapolation of that logic rather than a derived theorem about the infrared flow of $(8,8)$ SYM on arbitrary curved backgrounds.

The simplest dictionary is:
\begin{itemize}
\item The target Riemann surface $C$ of the first-order string is the physical spacetime on which the two-dimensional CFT is defined.
\item The eight free transverse bosons and fermions of the string are the same noncompact $\bR^8$ seed degrees of freedom that appear in each copy of the orbifold SCFT.
\item The closed-string coupling should scale as
\begin{equation}\label{eq:string-coupling-scaling}
g_s\sim \frac{1}{N},
\end{equation}
in the standard matrix-string normalization, so the genus expansion reproduces the $1/N$ expansion of the large-$N$ CFT.
\item The free-orbifold point should correspond to the free first-order string, while string interactions should come from the same twist-field/joining-splitting data that appear in matrix string theory.
\end{itemize}

This is meant to be a close two-dimensional analogue of the logic emphasized recently by Komatsu and Maity \cite{komatsu2025stringduals2dym}: the string is written directly in first-order variables adapted to the field-theory spacetime, and its perturbative coupling is identified with the large-$N$ expansion parameter of a two-dimensional gauge theory. The present proposal differs in two important ways. First, the target-space theory is maximally supersymmetric and carries the full $\bR^8$ transverse matter system. Second, the local longitudinal worldsheet theory is not merely topological; it is the finite-$\Aeff$ Gomis--Ooguri string, so the theory keeps a genuine string scale.

There is, however, a useful formal limit in which the longitudinal sector simplifies further. If one sends $\Aeff\to\infty$, then the deformation \eqref{eq:metric-deformation} drops out and the longitudinal theory reduces to the pure BF/Poisson-sigma core. Within the thermal genus-one dictionary derived below, this corresponds to $\mu=R/(2\Aeff)\to 0$, hence $p\to 1$ in the grand-canonical ensemble. This suggests a topological corner with unsuppressed long cycles. It is not the regime used in the thermal comparison below, and beyond that comparison I do not claim to have established a full dual interpretation of this limit.

\section[First Check: Genus One and Large-N Orbifold Free Energy]{First Check: Genus One and the Large-$N$ Orbifold Free Energy}
Take the target spacetime to be a Euclidean thermal torus with temporal period $\beta=T^{-1}$ and spatial circle of circumference $2\pi R$,
\begin{equation}\label{eq:thermal-target}
t_E\sim t_E+\beta,
\qquad
x\sim x+2\pi R.
\end{equation}
Here $T$ is the temperature and $R$ is the target-space spatial radius. Since the seed theory on $\bR^8$ is noncompact, the partition functions in this subsection are understood per unit transverse volume, equivalently with the center-of-mass zero modes divided out. The natural large-$N$ quantity to compare with the genus-one string free energy is then not the full canonical partition function of $\Sym^N(\bR^8)$, which contains pieces growing with $N$, but its stable order-$N^0$ single-cycle contribution. Equivalently, one may package this data in a grand-canonical generating function with fugacity $p$ for cycle length and then take the plethystic logarithm. This is the standard second-quantized string interpretation of symmetric-product orbifolds \cite{dijkgraaf1997ellipticgenerasymmetricproducts}. The orbifold manipulations in this subsection are derived directly in the note, while the worldsheet answer quoted later is imported from \cite{gomis2000nonrelativisticclosedstringtheory}. For the physical thermal spin structure $[\nu_t,\nu_x]=[1,0]$, the multiplicity matching will also be derived explicitly below.

Let $\mc H_1^{[\nu_t,\nu_x]}$ be the Hilbert space of one copy of the $\bR^8$ SCFT on a circle of circumference $2\pi R$, with spacetime fermion boundary conditions labeled by
\begin{equation}\label{eq:target-spin-structure}
\nu_t,\nu_x\in \{0,1\},
\end{equation}
where $\nu=0$ means periodic and $\nu=1$ means antiperiodic. In operator language, $\nu_x$ specifies the spatial-circle Hilbert space, while $\nu_t$ specifies whether the Euclidean-time trace is ordinary ($\nu_t=1$) or twisted by $(-1)^F$ ($\nu_t=0$). The thermal free energy corresponds to $\nu_t=1$. Write the corresponding seed trace as
\begin{equation}\label{eq:seed-trace}
Z_1^{[\nu_t,\nu_x]}(\beta,\theta)
=
\mathrm{Tr}_{\mc H_1^{[\nu_t,\nu_x]}}
\exp\!\left[
-\beta \frac{L_0+\bar L_0}{R}
+i\theta(L_0-\bar L_0)
\right]
=
\sum_{h,\bar h} D^{[\nu_t,\nu_x]}(h,\bar h)\,
e^{-\beta(h+\bar h)/R+i\theta(h-\bar h)}.
\end{equation}
Here $L_0$ and $\bar L_0$ are the left- and right-moving Virasoro zero modes of a single copy of the $\bR^8$ SCFT, $\theta$ is a bookkeeping angle for level matching, and $D^{[\nu_t,\nu_x]}(h,\bar h)$ counts seed states of conformal weights $(h,\bar h)$ in the chosen spacetime spin structure. I use $h$ and $\bar h$ to denote the actual numbers appearing in the Boltzmann factor, including the sector-dependent zero-point energies appropriate to the chosen NS/R and GSO conventions. Any overall vacuum normalization that is independent of the cycle fugacity $p$ is left implicit.

For the actual comparison to the Gomis--Ooguri thermal free energy, the physically relevant spatial spin structure is periodic,
\begin{equation}\label{eq:physical-spatial-spin}
\nu_x=0,
\end{equation}
because this is the supersymmetric boundary condition inherited from the 2d $(8,8)$ $U(N)$ SYM on a circle. In this case the $w$-cycle twisted sector has the standard matrix-string long-string description \cite{dijkgraaf1997matrixstring,arutyunov1998orbifoldscattering}. Let $c_{\rm seed}=12$ be the chiral central charge of one copy of the $\bR^8$ SCFT, and let $(\Delta,\bar\Delta)$ denote the plane conformal weights of a seed operator. Under the $w$-fold covering map
\begin{equation}\label{eq:covering-map}
z=t^w,
\end{equation}
the corresponding twisted-sector operator has weights
\begin{equation}\label{eq:twisted-weight-map}
\Delta_w-\frac{c_{\rm seed}}{24}
=
\frac{1}{w}\left(\Delta-\frac{c_{\rm seed}}{24}\right),
\qquad
\bar\Delta_w-\frac{c_{\rm seed}}{24}
=
\frac{1}{w}\left(\bar\Delta-\frac{c_{\rm seed}}{24}\right).
\end{equation}
Equivalently, the cylinder energies above the Ramond vacuum scale as $1/w$. With the standard Dijkgraaf--Verlinde--Verlinde lift of the cyclic permutation to the fermions, the periodic seed sector is Ramond and its vacuum already has $\Delta_0=\bar\Delta_0=c_{\rm seed}/24=1/2$, so the square-root Jacobian from the covering map is absorbed into the usual twisted-vacuum energy shift and the long-string Hilbert space is again in the Ramond sector rather than acquiring an extra parity flip for even $w$. An alternative even-$w$ lift would shift the operator-channel sector and would require a separate analysis. Throughout this note I adopt the standard DVV choice. Therefore the twisted-sector multiplicities relevant to the thermal comparison are in fact $w$-independent:
\begin{equation}\label{eq:w-independent-degeneracy}
D_w^{[1,0]}(h,\bar h)=D_{\rm RR}^{\mathrm{seed}}(h,\bar h).
\end{equation}
Here $D_{\rm RR}^{\mathrm{seed}}(h,\bar h)$ is the degeneracy of one copy of the $\bR^8$ SCFT with periodic spatial fermions and the ordinary thermal trace. The label ``RR'' refers to the fact that both left- and right-moving Hilbert spaces are in their Ramond sectors along the spatial circle; it does \emph{not} mean that the Euclidean-time boundary condition is also periodic.
It is convenient to package the same seed Hilbert-space data as the specialized trace
\begin{equation}\label{eq:seed-rr-trace}
Z_{\rm RR}^{\mathrm{seed}}(\beta,\theta)
:=
Z_1^{[1,0]}(\beta,\theta)
=
\sum_{h,\bar h} D_{\rm RR}^{\mathrm{seed}}(h,\bar h)\,
e^{-\beta(h+\bar h)/R+i\theta(h-\bar h)}.
\end{equation}
To make this trace fully explicit, define
\begin{equation}\label{eq:seed-thermal-modulus}
q=e^{-\beta/R+i\theta}=e^{2\pi i\tau_{\rm seed}},
\qquad
\bar q=e^{-\beta/R-i\theta}=e^{-2\pi i\bar\tau_{\rm seed}},
\qquad
\tau_{\rm seed}:=\frac{\theta+i\beta/R}{2\pi}.
\end{equation}
If $V_8$ denotes the volume of the noncompact transverse $\bR^8$, then the free seed trace per unit transverse volume is
\begin{equation}\label{eq:seed-rr-factorized}
\frac{Z_{\rm RR}^{\mathrm{seed}}(\beta,\theta)}{V_8}
=
\int_{\bR^8}\frac{d^8k}{(2\pi)^8}\,
e^{-\beta \Aeff k^2/(2R)}\,
\chi^{\mathrm{osc}}_{\rm RR}(q)\,
\bar\chi^{\mathrm{osc}}_{\rm RR}(\bar q).
\end{equation}
Evaluating the Gaussian integral gives
\begin{equation}\label{eq:seed-rr-gaussian}
\frac{Z_{\rm RR}^{\mathrm{seed}}(\beta,\theta)}{V_8}
=
\left(\frac{R}{2\pi \beta \Aeff}\right)^4
\chi^{\mathrm{osc}}_{\rm RR}(q)\,
\bar\chi^{\mathrm{osc}}_{\rm RR}(\bar q).
\end{equation}
Here $k\in \bR^8$ is the transverse momentum and
\begin{equation}\label{eq:seed-rr-oscillator}
\chi^{\mathrm{osc}}_{\rm RR}(q)
=
16\prod_{n=1}^\infty \frac{(1+q^n)^8}{(1-q^n)^8}
=
\frac{\vartheta_2(\tau_{\rm seed})^4}{\eta(\tau_{\rm seed})^{12}},
\end{equation}
with the same formula antiholomorphically for $\bar\chi^{\mathrm{osc}}_{\rm RR}(\bar q)$. The factor $16$ is the degeneracy of the eight left-moving Ramond zero modes; the right-moving sector contributes the same factor.
The first few oscillator coefficients are
\begin{equation}\label{eq:seed-rr-oscillator-expansion}
\chi^{\mathrm{osc}}_{\rm RR}(q)
=
16+256q+2304q^2+15360q^3+84224q^4+O(q^5).
\end{equation}
So the full seed trace begins with a $16\times 16=256$-state Ramond ground multiplet, followed by the expected tower of left- and right-moving oscillator excitations.
The same free transverse system can also be described in path-integral language, for one fixed diagonal spin structure of the free transverse matter system, by torus partition functions per unit transverse volume,
\begin{equation}\label{eq:transverse-path-integral-spin}
\frac{Z_{\perp}^{[\alpha,\beta]}(\tau,\bar\tau)}{V_8}
=
\frac{1}{(4\pi^2 \Aeff\,\mathrm{Im}\,\tau)^4}\,
\frac{\vartheta\!\begin{bmatrix}\alpha\\ \beta\end{bmatrix}\!(0|\tau)^4}{\eta(\tau)^{12}}
\frac{\overline{\vartheta\!\begin{bmatrix}\alpha\\ \beta\end{bmatrix}\!(0|\tau)^4}}{\bar\eta(\bar\tau)^{12}},
\qquad
\alpha,\beta\in\left\{0,\frac12\right\}.
\end{equation}
Here $\tau$ is the Euclidean worldsheet torus modulus, while $\alpha$ and $\beta$ label fermion boundary conditions around chosen torus cycles; this $\beta$ is a theta-characteristic label, not the inverse temperature appearing in \eqref{eq:thermal-target}. The formula \eqref{eq:transverse-path-integral-spin} is the partition function in one \emph{fixed} spin structure before any GSO sum. This notation should not be confused with the operator-channel character \eqref{eq:seed-rr-trace}: once one chooses which torus cycle is interpreted as Euclidean time, modular transformations can exchange the various theta characteristics. In the present genus-one comparison I work in the operator channel adapted to the long-string Hamiltonian \eqref{eq:long-string-hamiltonian}, so the explicit seed character \eqref{eq:seed-rr-trace}--\eqref{eq:seed-rr-oscillator-expansion} is the relevant object.

Denote by $\mc H_w$ the Hilbert space of a single length-$w$ twisted sector. A twisted sector built from a single cycle of length $w$ is a ``long string'' whose effective spatial circle has circumference $2\pi wR$. Its Hamiltonian is therefore
\begin{equation}\label{eq:long-string-hamiltonian}
H_w=\frac{L_0+\bar L_0}{wR}.
\end{equation}
Projecting onto the $\mathbb Z_w$-invariant subspace gives
\begin{equation}\label{eq:cyclic-projector}
P_w=\frac{1}{w}\sum_{n=0}^{w-1} g^n,
\qquad
g\cdot |h,\bar h\rangle
=
e^{2\pi i(h-\bar h)/w}|h,\bar h\rangle.
\end{equation}
Hence the thermal partition function of one cycle of length $w$, weighted by fugacity $p^w$, is
\begin{equation}\label{eq:one-cycle-trace}
\begin{aligned}
Z^{\mathrm{1\mbox{-}cycle}}_w(\beta,p)
&=
\mathrm{Tr}_{\mc H_w}\!\left(
P_w\,p^w e^{-\beta H_w}
\right)
\\
&=
\frac{p^w}{w}\sum_{n=0}^{w-1}\sum_{h,\bar h}D_w^{[\nu_t,\nu_x]}(h,\bar h)\,
\exp\!\left[
-\beta \frac{h+\bar h}{wR}
+2\pi i \frac{n}{w}(h-\bar h)
\right].
\end{aligned}
\end{equation}
The sum over $n$ imposes
\begin{equation}\label{eq:orbifold-level-matching}
h-\bar h\equiv 0 \pmod w,
\end{equation}
which is the usual long-string level-matching condition.
In the physical thermal spin structure, \eqref{eq:w-independent-degeneracy} and \eqref{eq:seed-rr-trace} turn \eqref{eq:one-cycle-trace} into the compact character formula
\begin{equation}\label{eq:one-cycle-trace-thermal}
Z^{\mathrm{1\mbox{-}cycle}}_w(\beta,p)\Big|_{[1,0]}
=
\frac{p^w}{w}\sum_{n=0}^{w-1}
Z_{\rm RR}^{\mathrm{seed}}\!\left(\frac{\beta}{w},\frac{2\pi n}{w}\right).
\end{equation}

The order-$N^0$ free energy is the standard plethystic free-gas answer built from these single-cycle states:
\begin{equation}\label{eq:orbifold-free-energy-sum}
\begin{aligned}
F^{(0)}_{\Sym}(\beta,p)
&=
-T\sum_{m=1}^\infty \frac{1}{m}
\sum_{w=1}^\infty Z^{\mathrm{1\mbox{-}cycle}}_w(m\beta,p^m)
\\
&=
-T\sum_{m=1}^\infty \sum_{w=1}^\infty \sum_{n=0}^{w-1}
\frac{p^{mw}}{wm}\sum_{h,\bar h}D_w^{[\nu_t,\nu_x]}(h,\bar h)\,
\exp\!\left[
-m\beta \frac{h+\bar h}{wR}
+2\pi i \frac{n}{w}(h-\bar h)
\right].
\end{aligned}
\end{equation}
Using $\log(1-x)=-\sum_{m\ge1}x^m/m$, this becomes
\begin{equation}\label{eq:orbifold-free-energy-log}
F^{(0)}_{\Sym}(\beta,p)
=
T\sum_{w=1}^\infty
\sum_{\substack{h,\bar h\\ h-\bar h\in w\mathbb Z}}
D_w^{[\nu_t,\nu_x]}(h,\bar h)\,
\log\!\left(
1-p^w e^{-\beta(h+\bar h)/(wR)}
\right).
\end{equation}
Equation \eqref{eq:orbifold-free-energy-log} is the derived orbifold answer in the present note. The universal $\log(1-x)$ is the standard symmetric-orbifold plethystic logarithm for cycle sectors. The spacetime fermion boundary conditions change the sector-dependent multiplicities $D_w^{[\nu_t,\nu_x]}(h,\bar h)$, not the combinatorial cycle factor itself.

In the physical thermal spin structure, \eqref{eq:one-cycle-trace-thermal} gives the more explicit character-level formula
\begin{equation}\label{eq:orbifold-free-energy-character}
F^{(0)}_{\Sym}(\beta,p)\Big|_{[1,0]}
=
-T\sum_{m=1}^\infty \sum_{w=1}^\infty \frac{p^{mw}}{wm}
\sum_{n=0}^{w-1}
Z_{\rm RR}^{\mathrm{seed}}\!\left(\frac{m\beta}{w},\frac{2\pi n}{w}\right),
\end{equation}
and equivalently, substituting \eqref{eq:w-independent-degeneracy} into \eqref{eq:orbifold-free-energy-log} gives the explicit logarithmic formula
\begin{equation}\label{eq:orbifold-free-energy-thermal}
F^{(0)}_{\Sym}(\beta,p)\Big|_{[1,0]}
=
T\sum_{w=1}^\infty
\sum_{\substack{h,\bar h\\ h-\bar h\in w\mathbb Z}}
D_{\rm RR}^{\mathrm{seed}}(h,\bar h)\,
\log\!\left(
1-p^w e^{-\beta(h+\bar h)/(wR)}
\right).
\end{equation}

Now compare with the genus-one worldsheet result of Gomis and Ooguri \cite{gomis2000nonrelativisticclosedstringtheory}. The following formula is quoted from their one-loop free-energy computation, rewritten in the present notation. Their torus path integral localizes to holomorphic maps from the worldsheet torus to the target thermal torus and gives
\begin{equation}\label{eq:go-free-energy-sum}
\begin{aligned}
F_{\mathrm{GO}}(T)
&=
-\sum_{m=1}^\infty \sum_{w=1}^\infty \sum_{n}
\frac{T}{wm}\sum_{h,\bar h}D^{\mathrm{GO}}_w(h,\bar h)\,
\exp\!\left[
-\frac{m}{T}\left(
\frac{wR}{2\Aeff}+\frac{h+\bar h}{wR}
\right)
+2\pi i \frac{n}{w}(h-\bar h)
\right],
\end{aligned}
\end{equation}
which they resum as
\begin{equation}\label{eq:go-free-energy-log}
F_{\mathrm{GO}}(T)
=
T\sum_{w=1}^\infty
\sum_{\substack{h,\bar h\\ h-\bar h\in w\mathbb Z}}
D^{\mathrm{GO}}_w(h,\bar h)\,
\log\!\left(
1-e^{-E(w,h,\bar h)/T}
\right),
\end{equation}
with
\begin{equation}\label{eq:go-energy}
E(w,h,\bar h)
=
\frac{wR}{2\Aeff}+\frac{h+\bar h}{wR}.
\end{equation}
In the notation of \cite{gomis2000nonrelativisticclosedstringtheory}, the quantities $h$ and $\bar h$ already include both transverse momenta and transverse oscillator excitations:
\begin{equation}\label{eq:go-transverse-weights}
h=\frac{\Aeff k^2}{4}+N,
\qquad
\bar h=\frac{\Aeff k^2}{4}+\bar N.
\end{equation}
Here $k\in \bR^8$ is the transverse momentum, while $N$ and $\bar N$ are the left- and right-moving transverse oscillator levels. Thus the multiplicity factor $D_w^{\mathrm{GO}}(h,\bar h)$ appearing in the Gomis--Ooguri answer is simply their transverse multiplicity $D(h,\bar h)$, already independent of $w$, written in the same seed Hilbert-space variables used above. For the thermal comparison at hand, I identify this multiplicity with the same transverse free $\bR^8$ lightcone Hilbert space used in \eqref{eq:seed-rr-trace}: eight bosons together with eight fermions in the Ramond spatial sector, evaluated in the ordinary thermal trace. The finite-temperature GSO choice itself is part of the quoted Gomis--Ooguri result; the present point is only that, after that choice is made, the surviving multiplicity is the same seed degeneracy. Thus
\begin{equation}\label{eq:go-degeneracy-identification}
D_w^{\mathrm{GO}}(h,\bar h)=D_{\rm RR}^{\mathrm{seed}}(h,\bar h).
\end{equation}
Because the phase in \eqref{eq:go-free-energy-sum} depends only on $n$ modulo $w$, one may choose representatives $n=0,\dots,w-1$. Using \eqref{eq:seed-rr-trace} and \eqref{eq:go-degeneracy-identification}, the quoted Gomis--Ooguri formula can then be rewritten as
\begin{equation}\label{eq:go-free-energy-character}
F_{\mathrm{GO}}(T)
=
-T\sum_{m=1}^\infty \sum_{w=1}^\infty \frac{e^{-m\beta wR/(2\Aeff)}}{wm}
\sum_{n=0}^{w-1}
Z_{\rm RR}^{\mathrm{seed}}\!\left(\frac{m\beta}{w},\frac{2\pi n}{w}\right),
\end{equation}
where $\beta=T^{-1}$.
If desired, one may now substitute the explicit free-field expression \eqref{eq:seed-rr-factorized}--\eqref{eq:seed-rr-oscillator} into either \eqref{eq:orbifold-free-energy-character} or \eqref{eq:go-free-energy-character}. This produces exactly the same momentum integral over $k$ and the same Ramond oscillator product in both formulas.
One can also repackage the same statement in the path-integral notation \eqref{eq:transverse-path-integral-spin}. To do so uniformly for arbitrary target-torus spin structures, however, one must track carefully the induced theta characteristic of the holomorphic map together with the modular transformation back to the operator channel used here. I have not attempted that more global bookkeeping in the present note, because for the physical thermal sector $[\nu_t,\nu_x]=[1,0]$ the operator-channel argument already suffices.
The algebraic dependence on $w$, $(h,\bar h)$, and the Boltzmann factor then matches the orbifold answer after the identification
\begin{equation}\label{eq:chemical-potential}
p=e^{-\beta\mu},
\qquad
\mu=\frac{R}{2\Aeff}.
\end{equation}
Here $\mu$ is the chemical potential conjugate to cycle length, equivalently to target winding in the first-order-string interpretation. In the physical thermal spin structure $[\nu_t,\nu_x]=[1,0]$, and with the standard DVV lift adopted above, the covering-map argument gives
\begin{equation}\label{eq:degeneracy-match}
D_w^{[1,0]}(h,\bar h)=D^{\mathrm{GO}}_w(h,\bar h),
\end{equation}
because both sides are equal to the same seed Ramond-sector degeneracy $D_{\rm RR}^{\mathrm{seed}}$ by \eqref{eq:w-independent-degeneracy} and \eqref{eq:go-degeneracy-identification}. Thus the multiplicity matching required for the thermal test is established within that standard long-string convention. Comparing \eqref{eq:orbifold-free-energy-character} with \eqref{eq:go-free-energy-character}, one finds a term-by-term equality once \eqref{eq:chemical-potential} is imposed. Equivalently, after substituting \eqref{eq:chemical-potential} into \eqref{eq:orbifold-free-energy-thermal} the orbifold free energy becomes
\begin{equation}\label{eq:free-energy-match}
F^{(0)}_{\Sym}(\beta,\mu)\Big|_{[1,0]}
=
T\sum_{w=1}^\infty
\sum_{\substack{h,\bar h\\ h-\bar h\in w\mathbb Z}}
D_{\rm RR}^{\mathrm{seed}}(h,\bar h)\,
\log\!\left(
1-e^{-\beta[\mu w +(h+\bar h)/(wR)]}
\right),
\end{equation}
which then reproduces the Gomis--Ooguri answer.

This comparison is stronger than a mere agreement of final spectra. The two auxiliary sums appearing in \cite{gomis2000nonrelativisticclosedstringtheory} acquire a direct orbifold meaning:
\begin{itemize}
\item the $n$-sum is precisely the $\mathbb Z_w$ projector onto cyclically invariant states in the length-$w$ twisted sector,
\item the $m$-sum is precisely the plethystic expansion of $\log(1-x)$.
\end{itemize}
So the worldsheet localization on holomorphic maps to the target torus found by Gomis and Ooguri is mirrored, on the orbifold side, by the standard cycle decomposition of the symmetric product. In the thermal spin structure $[\nu_t,\nu_x]=[1,0]$, and within the standard DVV lift, the match is therefore numerical as well as structural.

It is worth being precise about the extent of this nontriviality. The genuinely nontrivial part of the test is the longitudinal/stringy reorganization: the holomorphic-map sum appearing in the quoted genus-one Gomis--Ooguri formula, including the linear winding energy $wR/(2\Aeff)$, the $\mathbb Z_w$ level-matching projector, and the $m$-sum over plethystic images, is reproduced exactly by the large-$N$ single-cycle orbifold free energy after the identification \eqref{eq:chemical-potential}. That is a real constraint on the proposed duality.

At the same time, this genus-one/$N^0$ test is not yet decisive. Both sides use the same free transverse $\bR^8$ matter system, so once the standard DVV long-string lift is chosen, the detailed coefficient matching is still largely kinematical. In other words, the test shows that the first-order string reorganizes the free orbifold data in the correct string-theoretic way, but it does not yet probe genuinely interacting joining/splitting data or other observables that would distinguish the proposal more sharply.

There is still an additional subtlety for \emph{other} target-torus spin structures: if one wants a uniform formula for arbitrary $[\nu_t,\nu_x]$, then one must track carefully how the covering map acts on fermions in NS sectors and how $(-1)^F$ insertions are represented in the orbifold trace. That more general bookkeeping remains worthwhile, but it is no longer needed for the concrete thermal test carried out here.

There is one conceptual caveat. The exact formula above is naturally grand canonical in cycle length, because the linear piece $wR/(2\Aeff)$ is represented by the fugacity $p^w$. If one insists on the strictly canonical fixed-$N$ partition function, one should extract the coefficient of $p^N$ at the end. The large-$N$ order-$N^0$ data are nevertheless controlled by the same single-cycle building blocks written above, while finite-$N$ corrections and genuine string interactions should begin to deform this free result.

\section{Immediate Checks}
The next steps seem concrete.
\begin{itemize}
\item Formulate the BRST complex and GSO projection globally on curved $C$, using the curved super-$\beta\gamma$ patching above rather than the old second-order variables.
\item Extend the genus-one check above from the thermal torus to more general background Riemann surfaces and identify what replaces the long-cycle projector.
\item Show explicitly whether changing the target metric representative $e^{2\omega}$ is BRST trivial or at worst changes the theory by local counterterms with no effect on physical observables.
\item Identify the precise twist interaction that reconstructs genus expansion and joining/splitting vertices from the symmetric-orbifold side, in the same way matrix string theory reconstructs ordinary type-II string interactions.
\item Compare the first nontrivial observables on both sides, for example protected partition functions or low-point correlators on $C=\bC$, $C=S^2$, and $C=T^2$.
\end{itemize}

In short, the point of departure is simple. Start from the exact first-order string of \cite{gomis2000nonrelativisticclosedstringtheory}, keep the eight free transverse multiplets, and regard the longitudinal pair $(X^+,X^-)$ as coordinates on a genuine target Riemann surface. Nekrasov's curved $\beta\gamma$ framework \cite{nekrasov2005curvedbetagammasystems} is the essential tool for making this coordinate transformation problem globally well defined. What is derived here is the classical patching structure, the classical holomorphic-map solution locus, the orbifold bookkeeping formula \eqref{eq:orbifold-free-energy-log}, and the thermal-spin-structure reformulation of the quoted Gomis--Ooguri genus-one result that yields the matching \eqref{eq:degeneracy-match}--\eqref{eq:free-energy-match} for $[\nu_t,\nu_x]=[1,0]$ within the standard DVV long-string convention and the operator-channel conventions stated in the genus-one section above. That genus-one check is genuinely nontrivial in how it reconstructs the stringy longitudinal organization of the free orbifold, but it still remains a free-field consistency test rather than an interaction-level derivation of the duality. What remains conjectural or incomplete is the quantum completion: normal ordering of the curved deformation \eqref{eq:metric-deformation}, the anomaly problem summarized around \eqref{eq:target-euler}, the global BRST/GSO implementation on curved $C$, and the extension of the genus-one comparison to general target curves and to other target-torus spin structures.

\printbibliography

\end{document}
